% ===================================================================
%  সারণি নং 9 — পাহাড়। — জল-নগরী। — বন। — শিকার।
%
%  Ceci n'est PAS une transcription : c'est une traduction, faite en
%  2026, en bengali d'aujourd'hui. Voir texto/bn/10-sarani-01.tex
%  pour la consigne, qui vaut pour les seize fichiers.
% ===================================================================

%%K t09-apar-1 apar
\VUsaut{4.0mm}
\VUtitre{15.0pt}{60}{সারণি নং 9}
\VUsaut{3.0mm}

%%K t09-tit-1 sub
\VUcentre{12.0pt}{0}{\textit{প্রথম দৃশ্য।}}
\VUsaut{3.0mm}

%%K t09-tit-2 sub
\VUpk{11.4pt}{40}{পাহাড়। --- জল-নগরী।}
\VUsaut{3.0mm}

%%K t09-c1-01-1 p
1. --- আমি এক সপ্তাহ ধরে প্রাৎসে আছি। পিরেনিজের অপূর্ব \VUgras{পর্বতমালার}
\textsuperscript{(1)} সামনে দাঁড়িয়ে আমার যে বিস্ময় হয়েছিল, তা আমি বলে বোঝাতে
পারব না; তার \VUgras{চূড়া-রেখা} \textsuperscript{(2)} নীল আকাশে কোনো বিশাল
ঝালরের মতো ফুটে ওঠে।

%%K t09-c1-02-1 p
2. --- আমি কল্পনাও করিনি যে পিরেনিজের \VUgras{শিখর} \textsuperscript{(3)} এত
\VUgras{উচ্চতায়} পৌঁছাতে পারে, আর আমি সেই চমৎকার \VUgras{হিমবাহ}
\textsuperscript{(5)} ও বরফে ঢাকা \VUgras{চূড়া} \textsuperscript{(4)} দেখে হতবাক
হয়ে গেলাম, যেগুলি থেকে এত ভয়ংকর \VUgras{তুষারধস} নামে।

%%K t09-tit-3 sub
\VUsaut{3.0mm}
\VUpk{10.2pt}{40}{পাহাড়ি দৃশ্য।}
\VUsaut{3.0mm}

%%K t09-c1-03-1 p
3. --- আসার পরদিনই আমি প্রাৎসের কাছে সেই পুরনো নিভে যাওয়া \VUgras{আগ্নেয়গিরি}
\textsuperscript{(6)} পর্যন্ত গেলাম। \VUgras{জ্বালামুখের} খালি অনুর্বর কিনারায়
এখনও সেই চিহ্নগুলি স্পষ্ট দেখা যায় যা সেই \VUgras{লাভা} রেখে গিয়েছিল, যা
কয়েক শতাব্দী আগে \VUgras{পাহাড়ের ঢাল} \textsuperscript{(7)} বেয়ে বয়েছিল।
একটি \VUgras{ঢালে} আমি সেই লাভার কয়েকটি টুকরো পেয়েও গেলাম।

%%K t09-c1-04-1 p
4. --- প্রাৎস সীমান্তে বসানো, আর যে \VUgras{দুর্গ} \textsuperscript{(8)} ফ্রান্স
ও স্পেনের মধ্যেকার \VUgras{গিরিপথ} \textsuperscript{(27)} পাহারা দেয়, তা রাইনের
ধারের সেই পুরনো গড়গুলির মতোই যা নিজেদের পাথরের উপর থেকে কোনো \VUgras{শিকারের}
তাকে বসে আছে বলে মনে হয়।

%%K t09-c1-05-1 p
5. --- একটি বিশাল \VUgras{ঈগল} \textsuperscript{(9)}, যার বাসা \VUgras{গড়}
\textsuperscript{(8)} থেকে দূরে নয়, প্রায়ই আমার নজর কাড়ে। এই
\VUgras{শিকারি পাখি}, নিজের বলবান \VUgras{ঠোঁট} নিয়ে, প্রায়ই নিজের বাচ্চাদের
জন্য কোনো মেষশাবক বা ছাগলছানা নিজের \VUgras{থাবায়} চেপে আনে। তার
\VUgras{ডানা} এত বড় যে সে নিজের ভারী বোঝা নিয়ে দীর্ঘক্ষণ ভাসতে পারে।

%%K t09-tit-4 sub
\VUsaut{3.0mm}
\VUpk{10.2pt}{40}{পথিক।}
\VUsaut{3.0mm}

%%K t09-c1-06-1 p
6. --- এখানকার সবচেয়ে সুখকর বেড়ানো হল কো-র \VUgras{ঝরনা} \textsuperscript{(10)}
পর্যন্ত যাওয়া। \VUgras{পতিত জল} ঘূর্ণি তুলে নামে আর তারপর নগরের পশ্চিমের
\VUgras{দেবদারু বনে} \textsuperscript{(11)} একটি সুন্দর \VUgras{নালা} হয়ে
মিলিয়ে যায়। আমরা এই বনের ভিতর দিয়ে \VUgras{আবহাওয়া-কেন্দ্র}
\textsuperscript{(12)} পর্যন্ত উঠলাম, আর \VUgras{মিনারের} উপর থেকে অঞ্চলের সমস্ত
\VUgras{পাখির চোখে দেখা দৃশ্য} গুটিয়ে নিলাম। \VUgras{প্রাকৃতিক দৃশ্য}
অতুলনীয়, আর মনে হয় \VUgras{প্রকৃতি} এই দেশের উপরে নিজের সব ভান্ডার ঢেলে
দিয়েছে।

%%K t09-c1-07-1 p
7. --- নামার জন্য আমরা \VUgras{দড়ি-রেল} \textsuperscript{(13)} নিলাম আর একটি
ছোট \VUgras{স্টেশনে} থামলাম, একটি পুরনো \VUgras{সামন্ত দুর্গের}
\VUgras{ধ্বংসাবশেষের} \textsuperscript{(14)} পাশে, যেখানে একসময় প্রাৎসের
সামন্তেরা থাকতেন।

%%K t09-c1-08-1 p
8. --- বেচারা দুর্গ! সে কী ভাবে যখন নিজের নিচে সেই ছিমছাম \VUgras{জল-নগরী}
\textsuperscript{(15)} দেখে, যা আজ একটি ফ্যাশনদুরস্ত \VUgras{তাপীয় আরোগ্যস্থল}?

%%K t09-c1-08-2 p
\VUgras{জলবায়ু} এত সুখকর ও \VUgras{খনিজ জল} এত গুণকর যে
\VUgras{আরোগ্যশালায়} \textsuperscript{(17)} নেওয়া \VUgras{চিকিৎসাক্রম} সত্যিই
সুখের কথা।

%%K t09-tit-5 sub
\VUsaut{3.0mm}
\VUpk{10.2pt}{40}{একটি মনোরম জল-নগরী।}
\VUsaut{3.0mm}

%%K t09-c1-09-1 p
9. --- সত্যিই, থাকাটিকে সুখকর করতে সব কিছু করা হয়েছে। রেল নিয়ে যাওয়ার জন্য
একটি বড় \VUgras{সেতু-পুল} \textsuperscript{(16)} বানানো হয়েছে;
\VUgras{জলভবনের} \VUgras{উদ্যান} \textsuperscript{(18)}, যেখানে
\VUgras{সংগীতসভা} হয়, বড় সৌকর্যে সাজানো। কয়েকটি সুঠাম \VUgras{কুঠি}
\textsuperscript{(19)} \VUgras{জুয়াঘরের} \textsuperscript{(21)} চারপাশে গড়া, আর
খুব ভালোভাবে রাখা একটি \VUgras{রাজপথ} \textsuperscript{(22)} যাতায়াত খুব সহজ
করে দেয়।

%%K t09-c1-10-1 p
10. --- একটি সাদামাটা \VUgras{বাঁধ} \textsuperscript{(23)} \VUgras{রাস্তাকে}
\textsuperscript{(20)} আসল \VUgras{উপত্যকা} \textsuperscript{(25)} থেকে আলাদা করে,
যার তলদেশে একটি সুঠাম ছোট \VUgras{হ্রদ} \textsuperscript{(26)}, যা একটি স্বচ্ছ
\VUgras{নালাকে} \textsuperscript{(24)} জল জোগায়।

%%K t09-c1-11-1 p
11. --- উপত্যকার অন্য দিকে লোহার একটি \VUgras{খনি} \textsuperscript{(28)} খোঁড়া
হচ্ছে। কয়েকবার আমি \VUgras{খনিশ্রমিকদের} \VUgras{কুয়োয়} নামতে দেখতে গেছি,
আর আমি সেই \VUgras{সুড়ঙ্গেও} গেছি যেখানে তারা সারা দিন সেই \VUgras{আকরিক}
তোলে যাতে \VUgras{ধাতু} থাকে।

%%K t09-c1-12-1 p
12. --- খনির পাশে একটি বড় \VUgras{ঢালাইশালা} বানানো হয়েছে, যাতে তোলা সম্পদ
সেখানেই কাজে লাগে। \VUgras{চুল্লি} \textsuperscript{(29)}, যা এখানে প্রচুর
পাওয়া \VUgras{কয়লা} দিয়ে গরম করা হয়, তাড়াতাড়ি ও সস্তায় \VUgras{কাঁচা
লোহা} পেতে দেয়।

%%K t09-c1-13-1 p
13. --- খনি থেকে দূরে নয় একটি \VUgras{পাথরখাদ} \textsuperscript{(30)} কাটা
হচ্ছে। বুড়ো \VUgras{খাদ-মজুর} নিজের \VUgras{কোদাল} দিয়ে যা ভাঙে তা
\VUgras{গ্রানাইটও} নয়, \VUgras{পোরফিরিও} নয়, \VUgras{বেলেপাথরও} নয়, বরং
বাড়ি বানানোর সাধারণ \VUgras{ইমারতি পাথর}।

%%K t09-c1-14-1 p
14. --- এই দেশের সবচেয়ে সুন্দর অলংকারগুলির একটি \VUgras{দেবদারু}
\textsuperscript{(31)}। সর্বত্র --- কোনো \VUgras{খাদের} \textsuperscript{(32)}
তলায়, কোনো \VUgras{প্রপাতের} \textsuperscript{(33)} ধারে, কোনো \VUgras{গহ্বর}
\textsuperscript{(34)} বা খাদের পাশে --- তার সরু সুচগুলি নিজের সবুজ সুর জুড়ে
দেয়।

%%K t09-tit-6 sub
\VUsaut{3.0mm}
\VUpk{10.2pt}{40}{পায়ে হেঁটে চলা।}
\VUsaut{3.0mm}

%%K t09-c1-15-1 p
15. --- আমি নিজের ছুটির সদ্ব্যবহার করছি লম্বা বেড়ানোয়। যে
\VUgras{পায়ে-চলা পথগুলি} \textsuperscript{(35)} পাহাড়ে পাক খেয়ে ওঠে, তাতে
কয়েক \VUgras{কিলোমিটার}, আর প্রায়ই কয়েক \VUgras{ক্রোশ}, দূরত্ব টেরই না
পেয়ে পেরিয়ে যায়।

%%K t09-c1-15-2 p
\VUgras{মাইলফলক} \textsuperscript{(36)} ও \VUgras{দিকনির্দেশক}
\textsuperscript{(37)} পথগুলি চিহ্নিত করে, যেখানে প্রায়ই কোনো তরুণ
\VUgras{পাহাড়ি ছেলে} \textsuperscript{(38)} মেলে, বগলে \VUgras{ঝোলা}
\textsuperscript{(40)} নিয়ে ও একটি \VUgras{মারমট} \textsuperscript{(39)} তুলে,
বা কোনো \VUgras{যাযাবর} \textsuperscript{(41)}, যার \VUgras{লোমশ টুপি}
\textsuperscript{(42)} তার ছেঁড়া পুরনো \VUgras{ওড়নার} \textsuperscript{(43)}
সঙ্গে অদ্ভুত মিল খায়। সে একটি সধানো \VUgras{ভালুককে} \textsuperscript{(44)}
\VUgras{শিকল} ধরে নিয়ে চলে।

%%K t09-tit-7 sub
\VUpk{10.2pt}{40}{চড়াই।}
\VUsaut{3.0mm}

%%K t09-c1-16-1 p
16. --- প্রায় প্রতিদিন আমি একজন \VUgras{পথপ্রদর্শকের} \textsuperscript{(48)}
সঙ্গে \VUgras{আরোহণে} যাই, আর আমার বড় গর্ব হয় যখন পিঠে নিজের
\VUgras{ঝোলা} \textsuperscript{(47)} ও হাতে নিজের \VUgras{লাঠি} নিয়ে, কোনো
সত্যিকারের \VUgras{পর্যটকের} \textsuperscript{(46)} মতো, আমি \VUgras{পাথরে}
\textsuperscript{(45)} চড়াই করি।

%%K t09-c1-17-1 p
17. --- পাহাড়ের চূড়া থেকে সমতলের লোকেরা পিঁপড়ের চেয়ে বড় দেখায় না;
\VUgras{ভেড়ার পাল} \textsuperscript{(50)}, যা \VUgras{চারণভূমিতে}
\textsuperscript{(49)} চরে, ছেলেমেয়েদের খেলনা মনে হয়। একটি \VUgras{পাইন}
\textsuperscript{(51)}, বা \VUgras{কাঠের পাহাড়ি ঘর} \textsuperscript{(52)}
পর্যন্ত, সবুজের উপরে একটি বিন্দুর বেশি নয়।

%%K t09-c1-18-1 p
18. --- এই বেড়ানোয় আমাদের \VUgras{পায়ে-চলা পথে} \textsuperscript{(53)}
প্রায়ই কোনো \VUgras{শামোয়া-শিকারি} \textsuperscript{(54)} মেলে, যে এইমাত্র মারা
জন্তুটি কাঁধে নিয়ে থাকে।

%%K t09-c1-19-1 p
19. --- আমি নিজের উদ্ভিদ-সংগ্রহে \VUgras{ফার্নের} \textsuperscript{(55)} একটি খুব
অদ্ভুত পাতা ও \VUgras{হিদারের} \textsuperscript{(57)} একটি সুন্দর গোলাপি ডাল
রেখেছি, যা আমি গত বেড়ানোয় একটি ছোট \VUgras{গুহার} \textsuperscript{(56)} মুখে
কুড়িয়েছিলাম।

%%K t09-tit-8 sub
\VUsaut{3.0mm}
\VUcentre{12.0pt}{0}{\textit{দ্বিতীয় দৃশ্য।}}
\VUsaut{3.0mm}

%%K t09-tit-9 sub
\VUpk{11.4pt}{40}{বন। --- শিকার।}
\VUsaut{3.0mm}

%%K t09-c2-01-1 p
1. --- কাল আমি বনে একটি আনন্দময় দিন কাটিয়েছি। সেখানে বাস করা জন্তু ও
\VUgras{পোকামাকড়ের} \textsuperscript{(5)} ব্যস্ত জীবন আমার কাছে খুব
কৌতূহলজনক লাগল।

%%K t09-c2-01-2 p
\VUgras{মাকড়সা} \textsuperscript{(1)}, যে দুই ডালের মাঝে নিজের \VUgras{জাল}
\textsuperscript{(2)} বুনে পাশ দিয়ে যাওয়া \VUgras{মাছি} \textsuperscript{(3)}
ধরে, \VUgras{শামুক} \textsuperscript{(4)}, যে নিজের খোল নিয়ে কষ্টে এগোয়, সেই
গুবরে যে নিজের শিকারের খোঁজে, সেই পরিশ্রমী পিঁপড়ে যে \VUgras{পিঁপড়ের ঢিবির}
\textsuperscript{(6)} পাশে কাজে লেগে --- এই সব ছোট্ট প্রাণী অসাধারণ উৎসাহে
আসা-যাওয়া ও কাজ করতে থাকল।

%%K t09-c2-02-1 p
2. --- একটি \VUgras{সাপ} \textsuperscript{(7)}, যাকে প্রথমে আমি \VUgras{গোখরো}
ভেবেছিলাম কিন্তু যা কেবল একটি নিরীহ \VUgras{ঢোঁড়া} ছিল, একটি গুঁড়ির নিচে
কুণ্ডলী পাকিয়ে ছিল, আর মাঝে মাঝে নিজের সরু মাথা বের করছিল, যা সব সময় ছোবল
দিতে তৈরি বলে মনে হত। আমার সামনে একটি বড় \VUgras{শামুক-স্লাগ}
\textsuperscript{(8)} ভারী চালে বুকে হাঁটছিল, সেই সুস্বাদু ছোট্ট \VUgras{বুনো
স্ট্রবেরির} \textsuperscript{(11)} একটি সাবাড় করার পরে, যা এখানে প্রচুর জন্মায়।

%%K t09-c2-03-1 p
3. --- মাটি \VUgras{বুনো ফুলে} বিছানো ছিল: \VUgras{প্রিমরোজ}
\textsuperscript{(9)}, \VUgras{পেরিউইঙ্কল} \textsuperscript{(10)} ও আরও অনেক গাছ,
যা আমি চিনতাম না, \VUgras{ঘাসের গোছার} \textsuperscript{(12)} মাঝে ফুটে ছিল।

%%K t09-c2-04-1 p
4. --- এই অঞ্চলে \VUgras{ট্রাফল} প্রায় ততটাই সাধারণ যতটা \VUgras{ছত্রাক}
\textsuperscript{(14)}, আর এক বনরক্ষকের \VUgras{শুয়োরছানা} \textsuperscript{(13)}
লোভে খোঁজে লেগে ছিল যদি কিছু মেলে।

%%K t09-tit-10 sub
\VUsaut{3.0mm}
\VUpk{10.2pt}{40}{চোরা শিকার।}
\VUsaut{3.0mm}

%%K t09-c2-05-1 p
5. --- আমি একটি দৃশ্যের \VUgras{শেষ} দেখলাম যা আমাকে খুব হাসাল। নিজের
\VUgras{বাসে কুকুরের} \textsuperscript{(15)} নাকের কল্যাণে \VUgras{বনরক্ষক}
\textsuperscript{(16)} এইমাত্র একজন \VUgras{চোরা-শিকারিকে} \textsuperscript{(22)}
সেই মুহূর্তে হাতেনাতে ধরেছিলেন যখন সে নিজের মারা \VUgras{শিকার}
\textsuperscript{(17)} তুলে নিয়ে যাওয়ার তোড়জোড়ে ছিল। ধরা পড়ার চেয়ে তার
নিজের শিকার ছেড়ে দেওয়াই ভালো মনে হল, তাই চোরা-শিকারি পালাল, আর
\VUgras{খরগোশ} \textsuperscript{(18)}, \VUgras{হরিণী} \textsuperscript{(19)},
\VUgras{হরিণ} \textsuperscript{(20)} ও \VUgras{বুনো শুয়োর} \textsuperscript{(21)}
ঘাসের উপরে পড়ে রইল, বনরক্ষকের আনন্দের জন্য।

%%K t09-c2-06-1 p
6. --- বনে গাছের যত রকম আছে, তা হতবাক করে দেয়। \VUgras{বীচ}
\textsuperscript{(23)} ও \VUgras{ভূর্জ} \textsuperscript{(26)}, \VUgras{পাইন},
যা থেকে \VUgras{রজন} \textsuperscript{(27)} মেলে, \VUgras{ওক}
\textsuperscript{(29)} আর আরও অনেকে নিজেদের ডাল মেলায়।

%%K t09-c2-06-2 p
উঁচু গাছের গুঁড়িতে জড়ানো \VUgras{আইভি} \textsuperscript{(24)} \VUgras{উঁচু
বনকে} \textsuperscript{(25)} একটি ঝলমলে সবুজ রং দেয়, যা সেই \VUgras{শ্যাওলার}
\textsuperscript{(31)} ফিকে কোমল ছায়ার সঙ্গে সুখকরভাবে মিশে যায় যাতে
\VUgras{সবুজ কাঠঠোকরা} \textsuperscript{(30)} পোকা খোঁজে।

%%K t09-tit-11 sub
\VUsaut{3.0mm}
\VUpk{10.2pt}{40}{বনের দৃশ্য।}
\VUsaut{3.0mm}

%%K t09-c2-07-1 p
7. --- পুরনো ওকগাছ আমাকে মুগ্ধ করে। তাদের বড় পাক-খাওয়া \VUgras{শিকড়}
\textsuperscript{(32)}, আধেক মাটির বাইরে, তাদের বল ও ওজের চমৎকার রূপ দেয়, আর
আমার অবাক লাগে যে \VUgras{মালিক} \textsuperscript{(35)} কীভাবে তাদের কাটিয়ে সেই
\VUgras{করাতের} \textsuperscript{(34)} হাতে সঁপে দিতে রাজি হন যা
\VUgras{কাঠুরের} \textsuperscript{(33)}। বেচারা \VUgras{গুঁড়িগুলি}
\textsuperscript{(36)}, নিজেদের \VUgras{ছাল} \textsuperscript{(37)} ছাড়ানো বা
\VUgras{কুড়ুলে} \textsuperscript{(38)} চেরা, যা \VUgras{দিনমজুরের}
\textsuperscript{(39)}, আমাকে দুঃখ দেয়।

%%K t09-c2-08-1 p
8. --- কাছেই এক বুড়ো \VUgras{কয়লা-প্রস্তুতকারক} \textsuperscript{(41)} নিজের
\VUgras{বেলচা} দিয়ে কয়লার একটি \VUgras{স্তূপ} \textsuperscript{(42)} সাজিয়ে
রেখেছিল, আর এক বুড়ি কাটা গাছের ডালপালা দিয়ে বানানো \VUgras{শুকনো কাঠের}
একটি \VUgras{আঁটি} \textsuperscript{(40)} নিয়ে যাচ্ছিল।

%%K t09-tit-12 sub
\VUsaut{3.0mm}
\VUpk{10.2pt}{40}{শিকার।}
\VUsaut{3.0mm}

%%K t09-c2-09-1 p
9. --- আমি \VUgras{বনরক্ষকের ঘরে} \textsuperscript{(46)} একটু জিরোতে যাচ্ছিলাম,
কিন্তু সেই ছোট \VUgras{পুকুরে} \textsuperscript{(43)} পৌঁছে, যাতে একটি সুন্দর
\VUgras{রাজহাঁস} \textsuperscript{(45)} সাঁতরাচ্ছিল, দেখলাম যে সাম্প্রতিক একটি
\VUgras{বন্যা} \textsuperscript{(44)} পথটিকে সত্যিকারের \VUgras{জলাভূমি} বানিয়ে
দিয়েছে। আমাকে ঘুরপথে যেতে হল, আর \VUgras{সিডার} \textsuperscript{(49)},
\VUgras{পপলার} \textsuperscript{(48)} ও \VUgras{এলমের} \textsuperscript{(47)} পাশ
দিয়ে যেতে যেতে, যা বনের কিনারায় দাঁড়িয়ে, দেখলাম একটি \VUgras{ঝোপ}
\textsuperscript{(54)} থেকে বেরিয়ে একটি চমৎকার \VUgras{হরিণ} \textsuperscript{(50)}
ছুটল, যার পিছনে একটি ক্লান্তিহীন \VUgras{শিকারি কুকুরের পাল}
\textsuperscript{(51)} লেগে ছিল, যাকে \VUgras{শিঙা} \textsuperscript{(53)}
উত্তেজিত করছিল, \VUgras{ঘোড়সওয়ার শিকারিদের} \textsuperscript{(52)}। বেচারা
জন্তুটি একেবারে ক্লান্ত হয়ে পড়েছিল। সে আমার কয়েক পা দূরেই মরতে মরতে পড়ে
গেল।

%%K t09-c2-10-1 p
10. --- বনরক্ষকের ছেলে, যাকে \VUgras{শিকারের} শব্দ টেনে বের করল, ঘর থেকে
বেরোল, আর আমরা দুজনে আবার বনে ফিরলাম। সে একটি পুরনো ওকের ডালে একটি
\VUgras{নীলকণ্ঠ} \textsuperscript{(56)} দেখেছিল। তার মনে হল তার বাসা
\VUgras{পাতার} \textsuperscript{(58)} মধ্যে লুকানো থাকবে, আর সে সেটি নিতে
চাইছিল।

%%K t09-c2-10-2 p
\VUgras{কাঠবিড়ালির} \textsuperscript{(61)} মতো ক্ষিপ্র, সে এক \VUgras{ডাল}
\textsuperscript{(55)} থেকে অন্য ডালে চড়তে থাকল, কিন্তু কিছু পেল না আর কেবল এক
বুড়ো \VUgras{কাককে} \textsuperscript{(60)} চমকে দিতে পারল, যে একটি
\VUgras{শাখায়} \textsuperscript{(57)} বসে ছিল।
\VUsaut{4.0mm}
\VUfilet{22mm}
